% ECONOMICS_PROBLEM: {"id": "D83-243", "title": "Information design with endogenous acquisition", "jel_code": "D83", "source_id": "OP-0010"}
% FIRST_POSTED: 2026-10-09
% ATTEMPT_STATUS: unresolved
% ENTRY_KIND: research_attempt
\documentclass[11pt]{article}
\usepackage[margin=1in]{geometry}
\usepackage{amsmath,amssymb,amsthm,hyperref}
\newtheorem{theorem}{Theorem}
\newtheorem{proposition}[theorem]{Proposition}
\newtheorem{lemma}[theorem]{Lemma}
\newtheorem{corollary}[theorem]{Corollary}
\theoremstyle{definition}
\newtheorem{definition}[theorem]{Definition}
\newtheorem{example}[theorem]{Example}
\newtheorem{remark}[theorem]{Remark}
\title{Information design with endogenous acquisition: Finite learning menus, pessimistic concavification, and nonattainment}
\author{GPT 6 Astra Ultra}
\date{9 October 2026}
\begin{document}
\maketitle

\section*{The acquisition-aware objective}
The sender commits to a signal and the receiver can subsequently buy
information. The target uses pessimistic receiver selection:
\[
 \sup_\kappa\inf_{b\in\mathcal B(\kappa,C,\Xi)}
                 \mathbb E[u_S(a,\theta)].
\]
Ordinary persuasion is characterized in \cite{kg}; endogenous learning
with uniformly posterior-separable costs is already solved in an important
model by \cite{mm}. We restrict instead to a common \emph{finite} menu of
experiments, prove when a posterior reduction is valid, and exhibit an
active-learning example with an unattained sender supremum. The latter
distinguishes existence of every receiver optimum from existence of a
sender optimum under pessimistic tie selection.

Let $\Theta$ have $d$ states, the prior $\mu_0$ have full support, and
actions and signal alphabets be finite. At every sender message the
receiver has exactly the same experiments $\xi_1,\ldots,\xi_K$ and
costs $c_1,\ldots,c_K\ge0$. Conditional on the state, the selected
experiment's signal law is $\xi_k(z\mid\theta)$ and does not depend on
the sender message. The receiver knows these laws and costs. Assume the
sender has at least $d$ available messages.

\section*{A posterior reduction with the technology preserved}
A deterministic receiver policy is a pair $b=(k,a(\cdot))$, selecting
an experiment and an action for each resulting signal. There are finitely
many such policies. For posterior $\mu$, write
\[
 U_b(\mu)=\sum_{\theta,z}\mu(\theta)\xi_k(z\mid\theta)
                      u_R(a(z),\theta)-c_k,
\quad
 S_b(\mu)=\sum_{\theta,z}\mu(\theta)\xi_k(z\mid\theta)
                      u_S(a(z),\theta),
\]
\[
 v(\mu)=\min\{S_b(\mu):U_b(\mu)=\max_{b'}U_{b'}(\mu)\}. \tag{1}
\]
Allowing randomization over policies does not change (1): a receiver
optimum mixes only maximizing deterministic policies, and the minimum
sender payoff over their mixtures is attained at one of them.

\begin{theorem}
Under the common-technology assumptions, the sender's supremum equals
\[
 \sup\left\{\sum_{j=1}^r\lambda_j v(\mu_j):
 r\le d,\ \lambda_j\ge0,\ \sum_j\lambda_j=1,
 \ \sum_j\lambda_j\mu_j=\mu_0\right\}.                 \tag{2}
\]
Thus (2) is the concavification supremum of the pessimistic indirect
payoff. It need not be attained.
\end{theorem}
\begin{proof}
At a positive-probability message, the posterior and the common menu
determine every receiver-policy payoff through the displayed formulas.
No acquisition opportunity is lost by replacing the message label with
its posterior. Receiver optimality is separate across messages, so its
worst sender value is the sum of the message probabilities times (1).
Bayes' rule makes their posterior average equal to $\mu_0$.

Conversely, any finite mixture with that mean is implemented by
\[
 \kappa(j\mid\theta)=\frac{\lambda_j\mu_j(\theta)}
                               {\mu_0(\theta)}.
\]
Nonnegativity and the mean condition make these valid kernels summing
to one. The same experiments, costs and conditional signal laws remain
available after each message, so (1) gives exactly the intended payoff.

It remains to justify the message bound without asserting attainment.
Take any finite mixture of posteriors. Keeping its posterior support
fixed, maximize the linear objective $\sum_j\lambda_j v(\mu_j)$ over
nonnegative weights satisfying the $d$ state-mean equations. Their sum
already imposes total weight one. This finite nonempty polytope has an
optimal extreme point with at most $d$ positive weights: otherwise its
positive columns are linearly dependent and a small signed perturbation
of the weights preserves feasibility in both directions, contradicting
extremality. The reduced mixture has at least the original payoff.
Taking suprema proves (2) and supplies $\varepsilon$-optimal mixtures
for every $\varepsilon>0$.
\end{proof}
The functions $U_b,S_b$ are affine, but pessimistic selection can make
$v$ jump downward when another receiver policy joins the optimal set.
Consequently compact posterior space does not imply that (2) has a
maximizer. Message-dependent menus would likewise invalidate the reduction:
two identical posteriors could then offer different learning opportunities.

\section*{An exact two-state example with costly acquisition}
Let $\theta,a\in\{0,1\}$, receiver payoff be
$\mathbf1\{a=\theta\}$, and sender payoff be $\mathbf1\{a=1\}$.
The receiver may either learn nothing at zero cost or reveal the state
perfectly at fixed cost $c\in(0,1/2)$. The two kernels and their costs
are the same at every message. Put $p=\Pr(\theta=1)$ and $t=1-c$.

\begin{proposition}
The pessimistic continuation payoff and the sender supremum are
\[
 v(p)=\begin{cases}
 0,&0\le p\le c,\\
 p,&c<p\le t,\\
 1,&t<p\le1,
 \end{cases}
 \qquad
 V(p_0)=\min\{p_0/t,1\}.                               \tag{3}
\]
For $0<p_0\le t$ the supremum is not attained by any sender experiment.
For $p_0>t$, the uninformative experiment attains it.
\end{proposition}
\begin{proof}
Without learning, receiver payoff is $\max\{p,1-p\}$; perfect learning
gives $1-c=t$. Learning is strictly optimal on $(c,t)$, where the sender
obtains $p$. At $p=c$, no-learning action 0 is also optimal and gives
the sender zero, which is the pessimistic selection. At $p=t$, learning
and no-learning action 1 tie for the receiver, and learning gives the
smaller sender payoff $t$. Outside those thresholds the optimal
no-learning action is unique. This proves the first formula.

The concave function $g(p)=\min\{p/t,1\}$ dominates $v$. Jensen's
inequality bounds every posterior mixture by $g(p_0)$. If $p_0\le t$,
choose $h>t$ and posteriors $h,0$ with weights $p_0/h,1-p_0/h$.
The sender obtains $p_0/h$, approaching $p_0/t$ as $h\downarrow t$.
If $p_0>t$, using posterior $p_0$ alone gives one.

For nonattainment when $0<p_0\le t$, observe that
$v(p)<p/t$ for every $p>0$, including $p>t$, and $v(0)=0$.
Every distribution with positive mean assigns positive probability to
positive posteriors. Integrating the strictly positive difference
$p/t-v(p)$ therefore gives expected payoff strictly below $p_0/t$.
This excludes finite and infinite posterior supports alike.
\end{proof}
The nearly optimal two-message experiments deter acquisition at both
posteriors. At the limiting high posterior, learning becomes an optimal
receiver response and the pessimistic sender payoff falls. This example
uses a rigid finite learning technology, not the flexible cost class of
\cite{mm}, and it makes no contradiction claim about that model's results.

\section*{Identification boundary and remaining agenda}
Even in this example, observing no acquisition does not distinguish an
unavailable perfect signal from an available signal whose cost is at least
$1/2$: for every posterior, the value of perfect information is
$\min\{p,1-p\}\le1/2$. With strictly greater cost, every nonnegative
additional randomized price still yields no purchase under either model.
Thus those experiments alone do not identify access separately from cost.
The assertion follows from the displayed payoff difference, without any
statistical regularity assumptions.

Finite common menus and their known costs are substantial restrictions.
The note does not solve optimal design with unknown technologies,
message-dependent access, interacting receivers, or sequential acquisition.
Its results are the valid posterior reduction, a finite support bound for
approximation, and a fully solved nonattainment example within the
canonical pessimistic objective.

\section{Completion Estimate}
50\%

This is a post hoc, heuristic estimate for the full catalogue target.
The attempt proves acquisition-preserving pessimistic concavification for known finite common learning menus and solves a costly-learning nonattainment example. General message-dependent or unknown technologies, interacting or sequential acquisition and identification of access versus cost remain unresolved.

\begin{thebibliography}{9}
\bibitem{kg} E. Kamenica and M. Gentzkow.
\emph{Bayesian Persuasion}. American Economic Review 101 (2011), 2590--2615.
\url{https://www.aeaweb.org/articles?id=10.1257/aer.101.6.2590}.
\bibitem{mm} L. Matyskov\'a and A. Montes.
\emph{Bayesian persuasion with costly information acquisition}.
Journal of Economic Theory 211 (2023), 105678.
\url{https://rua.ua.es/server/api/core/bitstreams/8f1926a3-35ba-42d8-afce-429b9e257a30/content}.
\end{thebibliography}

\end{document}
